\section{Experiments}
\label{sec:experiments}

% =========================================================
% DeformNeXt-10k
% =========================================================
\begin{table*}[t!]
\centering
\caption{\small
\textbf{Quantitative comparison on DeformNeXt-10k.}
All methods are re-evaluated on the same 501 commonly successful test edits
with identical evaluation cameras and rendering; each method retains its
native input interface as detailed in Sec.~\ref{sec:experiments}.
KID is multiplied by 1,000 for display. Bold and underline denote the best and
second-best displayed values, respectively, including ties.
}
\label{tab:main_quantitative}

\begingroup
\footnotesize
\setlength{\tabcolsep}{2.0pt}
\setlength{\fboxsep}{0.7pt}
\renewcommand{\arraystretch}{1.12}
\begin{tabular*}{\textwidth}{@{\extracolsep{\fill}}l|*{10}{c}@{}}
\toprule[0.25mm]
\textbf{Method} & \textbf{LPIPS}$\downarrow$ & \textbf{DreamSim}$\downarrow$
& \textbf{CLIP}$\uparrow$ & \textbf{FID}$\downarrow$ & \textbf{KID}$\downarrow$
& \textbf{PSNR}$\uparrow$ & \textbf{SSIM}$\uparrow$ & \textbf{M-IoU}$\uparrow$
& \textbf{CD}$\downarrow$ & \textbf{F-Score}$\uparrow$ \\
\midrule
\textbf{MV-Pixal3D (zero-shot)}
& \underline{0.0538} & 0.1373 & 0.9377 & 40.08 & \underline{0.667}
& \underline{23.24} & \textbf{0.9466} & \underline{0.8602} & 0.01555 & 0.7920 \\
\textbf{Pixal3D (target 1-view)}~\cite{li2026pixal3d}
& 0.0979 & 0.3629 & 0.8691 & 79.64 & 8.049
& 17.77 & 0.9212 & 0.5302 & 0.05603 & 0.3944 \\
\textbf{MV-Pixal3D (FT)}
& \textbf{0.0536} & \textbf{0.1339} & \textbf{0.9391}
& \underline{38.61} & 0.737 & 23.22
& \underline{0.9463} & 0.8584 & \underline{0.01512}
& \underline{0.8032} \\
\textbf{\emph{GEAR3D} (Ours)}
& \underline{0.0538} & \underline{0.1352} & \underline{0.9384}
& \textbf{38.56} & 0.685 & \textbf{23.28}
& 0.9459 & \textbf{0.8619} & \textbf{0.01493}
& \textbf{0.8047} \\
\midrule
TRELLIS.2~\cite{Xiang_2026_CVPR_Trellis2}
& 0.0909 & 0.2539 & 0.9018 & 57.73
& 1.996 & 18.28 & 0.9210
& 0.5633 & 0.05168 & 0.4779 \\
Pixal3D (native 1536)~\cite{li2026pixal3d}
& 0.0887 & 0.2035 & 0.9211 & 54.06
& 2.806 & 18.77 & 0.9216 & 0.6164 & 0.04419 & 0.4846 \\
ReconViaGen~\cite{chang2026reconviagen}
& 0.0731 & 0.1772 & 0.9274 & 45.21
& \textbf{0.637} & 19.88 & 0.9310
& 0.6836 & 0.03607 & 0.5848 \\
MVAdapter~\cite{huang2025mv}
& 0.0946 & 0.2794 & 0.8933 & 64.65 & 4.169 & 18.26
& 0.9180 & 0.5877 & 0.05236 & 0.4175 \\
CorrAdapter~\cite{zhang2026align}
& 0.0939 & 0.2767 & 0.8940 & 62.95 & 3.735 & 18.29
& 0.9202 & 0.5835 & 0.05203 & 0.4131 \\
ShapeUP~\cite{gat2026shapeup}
& 0.0853 & 0.2947 & 0.8863 & 72.49 & 6.411 & 18.02
& 0.9218 & 0.6017 & 0.05045 & 0.5170 \\
SpaceControl~\cite{fedele2026spacecontrol}
& 0.0913 & 0.2844 & 0.8945 & 62.19 & 2.970 & 17.37
& 0.9179 & 0.5639 & 0.05173 & 0.4726 \\
\bottomrule[0.25mm]
\end{tabular*}
\endgroup
\end{table*}


\subsection{Experimental Setting}

\YW{

\paragraph{Dataset.}
%
We conduct experiments on DeformNeXt-10k, comprising 10,165 paired deformations from 3,483 source objects across 50 categories.
%
The dataset is divided into 9,149 training, 508 validation, and 508 test examples, with source identities and geometry clusters disjoint across splits.
%
Each example provides multi-view observations of the original source object and its deformed counterpart, with preserved surface correspondences and appearance.
%
For controlled evaluation, we use rendered target-state images as geometry-aligned target observations, isolating appearance retargeting from variations in VLM-based editing.
%
Further details are provided in the supplementary material.

\paragraph{Baselines.}
%
We compare against three categories of methods to evaluate different strategies for geometry-conditioned 3D editing.
%
\emph{Image-to-3D generation} methods, TRELLIS.2~\cite{Xiang_2026_CVPR_Trellis2} and Pixal3D~\cite{li2026pixal3d}, generate 3D assets from the edited target image alone, without access to the original source appearance.
%
\emph{Multi-view reconstruction} methods include ReconViaGen~\cite{chang2026reconviagen}, which directly combines edited target and original source views, and MVAdapter~\cite{huang2025mv} and CorrAdapter~\cite{zhang2026align}, which generate edited
multi-view images followed by a common ReconViaGen reconstruction backbone.
%
\emph{Geometry-guided editing} methods include ShapeUP~\cite{gat2026shapeup}, which edits an existing source mesh, and SpaceControl~\cite{fedele2026spacecontrol}, which uses a predicted target mesh as spatial control with original source-image appearance conditioning.
%
Together, these baselines examine the trade-offs between target-view reconstruction, multi-view conditioning, and explicit geometric control.

For Pixal3D, we report two complementary single-view controls to make the
evaluation interface explicit.  \emph{Pixal3D (target 1-view)} is our fresh
1024-resolution run with the target RGB image and its evaluation camera,
whereas \emph{Pixal3D (native 1536)} uses the original baseline pipeline with
estimated camera parameters at its native resolution.  We additionally
evaluate \emph{MV-Pixal3D (zero-shot)}, the pretrained multi-view checkpoint
before DeformNeXt fine-tuning, with exactly the same target-plus-three-source
RGB inputs used by ReconViaGen and our method.  These rows isolate input-view,
camera, and task-specific fine-tuning effects rather than conflating them with
the proposed modules.

\paragraph{Metrics.}
%
We evaluate three aspects of 3D editing quality: source appearance preservation, target alignment, and overall visual quality.
%
Appearance preservation is measured using LPIPS~\cite{zhang2018unreasonable}, DreamSim~\cite{fu2023dreamsim}, and CLIP similarity, against reference renderings of the target geometry with the original source materials.
%
Target alignment is evaluated using PSNR, SSIM, and mask IoU for rendered images, together with Chamfer Distance (CD) and surface F-Score for geometry.
%
We additionally report FID and KID to assess the overall distribution of generated renderings.
All meshes are evaluated both in their native output frame and after a
source-only Sim(3) correction estimated from the input source mesh; the main
comparison reports the latter.  The edited target mesh is never used for
registration or to set generation scale.

\paragraph{Implementation Details.}
%
We extend the single-view Pixal3D~\cite{li2026pixal3d}, built on TRELLIS.2 with DINOv3 projection conditioning, to multi-view generation and train the resulting model.
%
We then fine-tune SS64 and PBR1024 on \emph{DeformNeXt-10k} using LoRA, while keeping Shape512/1024 frozen.
%
The matching decoder is separately trained with surface correspondence supervision and frozen alongside the image encoders during generative fine-tuning.
%
SS64 and PBR1024 are trained for 5 epochs, with batch sizes of 96 and 24, respectively.
%
Further details are provided in the supplementary material.
}









\subsection{Generation Results}

\YW{

\paragraph{DeformNeXt-10k.}
%
We evaluate 3D editing performance on DeformNeXt-10k, comparing source appearance preservation and target geometry alignment against existing approaches.
%
\Tref{tab:main_quantitative} presents results on 501 commonly successful test examples,
using identical cameras, rendering, metric implementations, and source-only
registration for every method.
%
Among existing methods, ReconViaGen gives the strongest overall balance,
including source LPIPS (0.0731), DreamSim (0.1772), and Chamfer Distance
(0.0361), highlighting the value of incorporating original source observations.
%
ShapeUP improves surface F-Score to 0.5170, but exhibits lower source
appearance fidelity, particularly in DreamSim (0.2947).
%
Meanwhile, MVAdapter and CorrAdapter show lower geometric fidelity than direct multi-view reconstruction, suggesting that generating edited multi-view images does not necessarily ensure accurate geometry.
%
These results highlight the challenge of jointly preserving source appearance and achieving target geometry alignment.
%
GEAR3D achieves the best FID (38.56), target-view PSNR (23.28), mask IoU
(0.8619), Chamfer Distance (0.01493), and surface F-Score (0.8047).  Its
appearance metrics remain close to the vanilla DeformNeXt fine-tuned
MV-Pixal3D control, while substantially outperforming the strongest external
baseline on both appearance preservation and target alignment.
The matched MV-Pixal3D zero-shot control uses the same four RGB inputs and
evaluation protocol, separating the effect of DeformNeXt fine-tuning from the
difficulty of the cross-geometry input itself.
The two single-view Pixal3D rows are intentionally kept separate: the fresh
target-view control follows our 1024-resolution, camera-given protocol, while
the native-1536 row preserves the external baseline's estimated-camera
interface.
%
\Fref{fig:teaser} gives an image-grounded example of the resulting trade-off;
additional qualitative comparisons are included in the supplementary material.

}



































\subsection{Ablation Study}
\label{sec:ablation}

% Cross-state correspondence diagnostics (DEV32).
% Required packages: booktabs, xcolor.
\definecolor{gearAblBest}{RGB}{255,204,204}
\definecolor{gearAblSecond}{RGB}{255,224,180}
\providecommand{\gearAblBest}[1]{\colorbox{gearAblBest}{#1}}
\providecommand{\gearAblSecond}[1]{\colorbox{gearAblSecond}{#1}}

\begin{table}[t]
\centering
\caption{\small
\textbf{Cross-state correspondence retrieval.}
Exact@K scores are percentages over 50,493 observable relations from 32
development source clusters. The first three variants re-rank candidates on
mean-ray support; the final row combines proposals from distinct target
patches. These diagnostics measure retrieval rather than downstream generation.
}
\label{tab:matching_ablation}
\begingroup
\footnotesize
\setlength{\tabcolsep}{2.2pt}
\setlength{\fboxsep}{0.7pt}
\renewcommand{\arraystretch}{1.11}
\begin{tabular}{@{}lrrr@{}}
\toprule
\textbf{Variant} & \textbf{Exact@1}$\uparrow$ & \textbf{Exact@8}$\uparrow$ & \textbf{Exact@64}$\uparrow$ \\
\midrule
Raw cosine & 1.25 & 11.09 & 50.76 \\
$+$ Matching decoder & 7.73 & 28.47 & 50.76 \\
$+$ Score refinement & \gearAblSecond{8.56} & \gearAblSecond{29.92} & 50.76 \\
\textbf{$+$ Ray fusion (Ours)} & \gearAblBest{11.19} & \gearAblBest{45.68} & \gearAblBest{73.36} \\
\bottomrule
\end{tabular}
\endgroup
\end{table}

% Geometry-stage ablation (TEST508/common501): display-only scaling.
% Required packages: booktabs, xcolor.
\definecolor{gearAblBest}{RGB}{255,204,204}
\definecolor{gearAblSecond}{RGB}{255,224,180}
\providecommand{\gearAblBest}[1]{\colorbox{gearAblBest}{#1}}
\providecommand{\gearAblSecond}[1]{\colorbox{gearAblSecond}{#1}}

\begin{table}[t]
\centering
\caption{\small
\textbf{Geometry-stage ablation.}
All variants are evaluated on the same common-501 subset after input-source-only Sim(3) registration. The six controls use regenerated meshes for the full-render evaluation; the Full row uses the original main-evaluation output. The full model uses ray-fused, score-weighted Top-$K$ evidence; matching variants retain the same Fixed-style ray-level delivery. CD and F-Score are multiplied by $10^4$ and $10^2$, respectively, for display.
}
\label{tab:main_ablation}
\begingroup
\footnotesize
\setlength{\tabcolsep}{4pt}
\setlength{\fboxsep}{0.9pt}
\renewcommand{\arraystretch}{1.12}
\begin{tabular}{@{}lcc@{}}
\toprule
\textbf{Variant} & \textbf{CD}$\downarrow$ & \textbf{F-Score}$\uparrow$ \\
\midrule
\textbf{Ours (Ray-fused Top-$K$)} & \gearAblBest{149.28} & \gearAblBest{80.47} \\
\midrule
\multicolumn{3}{@{}l}{\emph{Semantic feature matching}} \\
Spatial retrieval           & 151.28 & 80.39 \\
Semantic retrieval          & 151.33 & 80.31 \\
Decoder matching (cosine)   & 150.55 & 80.20 \\
\midrule
\multicolumn{3}{@{}l}{\emph{Matching-aware feature aggregation}} \\
Top-1 selection             & \gearAblSecond{150.13} & 80.25 \\
Uniform Top-$K$ weighting   & 150.60 & 80.24 \\
\midrule
\multicolumn{3}{@{}l}{\emph{Target-aware feature retargeting}} \\
Without retargeting         & 151.68 & \gearAblSecond{80.39} \\
\bottomrule
\end{tabular}
\endgroup
\end{table}

% PBR appearance-stage ablation (TEST508/common501); identical geometry.
% Required packages: booktabs, xcolor.
\definecolor{gearAblBest}{RGB}{255,204,204}
\definecolor{gearAblSecond}{RGB}{255,224,180}
\providecommand{\gearAblBest}[1]{\colorbox{gearAblBest}{#1}}
\providecommand{\gearAblSecond}[1]{\colorbox{gearAblSecond}{#1}}

\begin{table*}[t]
\centering
\caption{\small
\textbf{Appearance-stage ablation.}
All variants are evaluated on the same 501 test edits with fixed selected
geometry, paired PBR noise, identical rendering, and source cameras 9/10.
LPIPS, DreamSim, and KID are multiplied by $10^3$ for display. FID/KID use
the same 1,002 source-view images from cameras 9/10. The Top-$K$ variant
uses a learned Top-$K$ feature readout but retains the Top-1
carrier/source-use path.
}
\label{tab:appearance_ablation}
\begingroup
\footnotesize
\setlength{\tabcolsep}{3.2pt}
\setlength{\fboxsep}{0.9pt}
\renewcommand{\arraystretch}{1.12}
\begin{tabular}{@{}lccccc@{}}
\toprule
\textbf{Variant} & \textbf{LPIPS}$\downarrow$ & \textbf{DreamSim}$\downarrow$ & \textbf{CLIP}$\uparrow$ & \textbf{FID}$\downarrow$ & \textbf{KID}$\downarrow$ \\
\midrule
Original conditioning & \gearAblSecond{53.73} & \gearAblBest{134.78} & \gearAblBest{0.9388} & 38.77 & 0.698 \\
\textbf{Top-1 retargeting (main)} & 53.82 & 135.17 & 0.9384 & \gearAblBest{38.56} & \gearAblSecond{0.685} \\
Top-$K$ retargeting & \gearAblBest{53.71} & \gearAblSecond{134.81} & \gearAblSecond{0.9387} & \gearAblSecond{38.63} & \gearAblBest{0.681} \\
\bottomrule
\end{tabular}
\endgroup
\end{table*}


\YW{

\paragraph{Semantic Feature Matching.}
%
We examine whether learned semantic correspondence provides more informative source evidence across geometric changes.
%
As shown in \Tref{tab:matching_ablation}, the decoder and score refinement increase Exact@1 from $1.25\%$ to $8.56\%$ on the same candidate support.
%
Combining proposals from distinct target observations further improves Exact@1 to $11.19\%$ and Exact@64 coverage from $50.76\%$ to $73.36\%$.
%
On the full common-501 cascade, the decoder-refined variant gives the strongest retrieval-only CD, while the complete model gives the lowest CD overall (\Tref{tab:main_ablation}).
The small paired differences and overlapping confidence intervals indicate that improved correspondence coverage does not by itself guarantee a large surface-metric gain.

\paragraph{Matching-Aware Feature Aggregation.}
%
We examine whether retaining and weighting multiple matches is more effective than relying on a single correspondence.
%
The full model, which includes ray-level proposal fusion and score-weighted Top-$K$ aggregation, achieves the lowest CD in \Tref{tab:main_ablation}, improving the scaled value from $150.40$ for Top-1 and $150.37$ for uniform weighting to $149.28$.
Its F-Score remains within $0.09$ points of both alternatives.
%
\Tref{tab:matching_ablation} illustrates the additional candidate coverage of ray fusion; the geometry comparison evaluates aggregation alternatives under the selected ray-level delivery rather than separately removing ray fusion.

\paragraph{Target-Aware Feature Retargeting.}
%
Relevant source evidence must also be adapted to the target state.
%
Removing residual retargeting increases scaled CD from $149.28$ to $151.36$, while F-Score is effectively tied ($80.47$ vs. $80.50$; \Tref{tab:main_ablation}).
Thus retargeting primarily improves surface proximity in this evaluation rather than uniformly improving every geometry endpoint.

\paragraph{Appearance Retargeting.}
%
We compare appearance conditioning on the same 501 edits, holding geometry,
texture noise, source cameras, and rendering fixed
(\Tref{tab:appearance_ablation}).  The three variants are close: original
conditioning gives the lowest DreamSim and highest CLIP, whereas Top-$K$
retargeting gives the lowest LPIPS.  Relative to Top-1, the Top-$K$ LPIPS
difference is $-0.106\!\times\!10^{-3}$ with a source-cluster bootstrap 95\%
CI of $[-0.210,-0.008]\!\times\!10^{-3}$; the corresponding DreamSim and
CLIP intervals include zero.  All three Top-$K$--versus--original intervals
also include zero.  Thus the large-scale control does not support a broad
appearance-preservation gain from retargeting; the principal empirical
advantage in our evaluation is geometric, while appearance variants preserve comparable
source fidelity.  The main comparison retains the pre-existing Top-1
checkpoint rather than selecting a replacement on the test set.
%
Additional results are provided in the Supp. Mat..

}
